\documentclass[10pt]{article}
\usepackage[margin=0.75in]{geometry}
\usepackage{amsmath,amssymb,amsthm,microtype}
\newtheorem{theorem}{Theorem}
\begin{document}
\begin{center}
{\Large A Forward-Horizon Obstruction for the Current Merge Table}\par\medskip
Collatz proof project\quad---\quad September 5, 2026
\end{center}
The current table contains 360 verified rules of the form
$n_0+AQ\mapsto x_0+BQ$, for natural $Q$, with smaller positive outputs
and exact merging paths. This note studies only that particular table.
Write $T(n)=n/2$ for even $n$ and $(3n+1)/2$ for odd $n$.

\begin{theorem}[Finite-table congruence classification; Lean checked]
Put $L=2^{18}3^{11}=46438023168$. For $0\leq j\leq11$, let $c_j$
be the unique residue modulo $L$ satisfying
\[
c_j\equiv-1\pmod{2^{18}},\qquad
c_j\equiv3^j2^{-j}-1\pmod{3^{11}},
\]
where $2^{-j}$ denotes modular inversion. If a listed rule applies to
$n\equiv c_j\pmod L$, it must be the inverse odd rule
\[
2+3Q\longmapsto1+2Q.
\]
No listed rule applies when $n\equiv c_0\pmod L$.
\end{theorem}
\begin{proof}
Every rule step $A$ divides $L$. The finite check verifies, for every
non-inverse-odd rule and each of the twelve classes,
$c_j\not\equiv n_0\pmod A$. Any application would contradict this
check. The remaining rule requires residue two modulo three, whereas
$c_0$ is divisible by three.

The exact representatives are
\[
\begin{array}{rrrr}
42120773631&16743137279&25114705919&14453047295\\
21679570943&9300344831&13950517247&44144787455\\
42998169599&41278242815&15479341055&46438023167.
\end{array}
\]
Lean checks the exclusions for these explicit representatives by ordinary
\texttt{decide}, and proves their extension to every natural quotient.
The Chinese-remainder interpretation is reproduced by the exact Python
certificate; it is not an additional Lean theorem.
\end{proof}

\begin{theorem}[Arbitrary fixed horizon; written deduction]
For every natural $K$, there are arbitrarily large positive seeds $N$
with noncontracting prefixes through $K$ such that, for every $k\leq K$,
every chain of the listed rules starting at $T^k(N)$ stops at or above
$N$. Chain length is unrestricted. In particular, no such chain supplies
a smaller replacement for $N$.
\end{theorem}
\begin{proof}
Choose any positive solution of
\[
N\equiv-1\pmod{2^{K+18}},\qquad N\equiv0\pmod{3^{11}}.
\]
The moduli are coprime, so solutions form an unbounded arithmetic
progression. The first $K$ steps are odd, and
\[
y_k:=T^k(N)=\frac{3^k(N+1)}{2^k}-1\qquad(0\leq k\leq K).
\]
Consequently $y_k\equiv-1\pmod{2^{18}}$ and
$y_k\equiv3^k2^{-k}-1\pmod{3^{11}}$. For $k\geq11$ the latter
residue is stationary at $-1$. Thus $y_k$ belongs to class
$c_{\min(k,11)}$.

For $k>0$, the only possible rule is the inverse odd rule, and its output
is exactly $y_{k-1}$, since $3y_{k-1}+1=2y_k$. At $k=0$ there is no
listed rule. Induction therefore restricts every chain to
$y_k,y_{k-1},\ldots,N$. All these seeds are at least $N$, and the
prefix coefficients are $(3/2)^i\geq1$. Under deficit transport the
initial deficit $(3/2)^k$ reaches one only on returning to $N$.
\end{proof}

\paragraph{Formal status and reproduction.}
The first theorem's classification of the explicit residues is kernel
checked in \texttt{Collatz.MergeRuleObstruction}; both public theorems
enter the axiom audit. The arbitrary-horizon deduction is a written proof,
not yet formalized in Lean. \texttt{analysis/merge\_rule\_obstruction.py}
reconstructs the residues and the CRT seeds; independent tests replay
horizons through 128 and multiple progression quotients. Those examples
are checks of the construction, not the reason its quantifier over $K$
holds.

\paragraph{Implication for the proof program.}
No uniformly bounded forward horizon combined with arbitrarily long
chains from this fixed table covers all finite-prefix candidates.
This does not refute coverage with a seed-dependent horizon, a larger
rule family, or another strategy. The constructed seeds are not asserted
to be fully noncontracting or divergent. Collatz remains unproved;
no mathematical priority claim is made.
\end{document}
